\section{Recovering the local fields after regrouping}
\label{sec:ope-singletons}
The preceding regrouping gives a polynomial density with controlled
coefficients. To continue the scale induction, its linear coefficients
must again be presented as local fields, with the remaining terms
assigned to interactions involving at least two marks. We give this
passage explicitly. It uses no Taylor approximation: the separation
by the number of marks is exact.

Fix a finite periodic lattice, a finite mark set $V$, and positive
weights $w_v$. Write $e_A$ for the squarefree monomial belonging to
an allowed mark set $A$, as in the preceding section. Ordinary
polynomial activities $F_X$ are indexed by supports $X$; flagged
activities $K_{X,Z}$ are indexed by pairs with
$\varnothing\ne Z\subseteq X$. Their coefficients are measurable,
gauge invariant, and depend only on the field and flag data in $X$.
A nonzero coefficient with mark set $A$ requires
$\{o_v:v\in A\}\subseteq X$. Ordinary activities have zero constant
coefficient and vanish unless $X$ is healthy. Flagged activities vanish
unless the flags in $X$ are exactly $Z$.

Put $\nu=\nu_R$, $M=M_R$ for the support parameters of the current
representation. Suppose the ordinary and flagged arrays have the
weighted coefficient bounds from Section~\ref{sec:ope-grouping}, with
local hit bounds $a$ and $b$, respectively. In concrete terms there
are nonnegative coefficient majorants $t_{X,A}$ and $s_{(X,Z),A}$,
valid for every field configuration, such that for every site $x$
\begin{align}
 \sum_{X\ni x}e^{\nu c_M(X)}
       \sum_A t_{X,A}\prod_{v\in A}w_v&\le a,\label{eq:ope-single-hit}\\
 \sum_{(X,Z):x\in X}e^{\nu c_M(X)}
       \sum_A s_{(X,Z),A}\prod_{v\in A}w_v&\le b.\notag
\end{align}
The underlying lattice is finite, so the support indices in these
expressions are finite. The activity dictionary used to represent
them may nevertheless be countable.

We use the following precise property of that dictionary. Every
nonempty support has an unused ordinary slot with that support and
zero background coefficient. Every admissible flagged pair has a
flagged slot with that pair, eligible whenever the specified flag
condition holds. A real measurable gauge-invariant function $h$ in
such a slot is an allowed elementary insertion when it has the
required local dependence and vanishes off the slot's eligibility
set. If $|h|\le t$, its cost is
\[
                 e^{\nu c_M(X)}t.
\]
This is a property of the available representation, not a consequence
of the support count. Its production in the evolving activity
dictionary remains an input to this local argument.

\subsection{The price of a singleton}
Fix a mark $v$. In the first inequality of
\eqref{eq:ope-single-hit}, take $x=o_v$ and retain only the
coefficient of $e_{\{v\}}$. Coefficients on supports not containing
$o_v$ vanish. Dividing by $w_v>0$ gives majorants $t_X\ge0$ with
\[
 |[e_{\{v\}}]F_X(U)|\le t_X,
 \qquad \sum_Xe^{\nu c_M(X)}t_X\le a/w_v.
\]
The same argument for the flagged array gives majorants $s_{X,Z}$
with weighted sum at most $b/w_v$. There is no factor $|V|$ in
these bounds: only one coordinate of the price at $o_v$ was retained.

For an ordinary singleton function $f$ supported in $X$, use the
unused ordinary slot and write
$f=\operatorname{Re}f+\mathrm i\operatorname{Im}f$.
Both real functions have bound $t_X$, retain local dependence,
measurability and gauge invariance, and vanish outside the healthy
set. They are therefore two elementary insertions, with coefficients
$1$ and $\mathrm i$, of total cost at most
$2e^{\nu c_M(X)}t_X$. They contribute $f$ to the ordinary field and
zero to its flagged profile. The same construction in the flagged
slot contributes zero to the ordinary field and $f$ at the specified
flagged pair, at cost at most $2e^{\nu c_M(X)}s_{X,Z}$.
Zero functions require no slot.

Consequently, for arbitrary constants $u_v\in\mathbb C$, the pair
\begin{equation}\label{eq:ope-single-pair}
 g_v(U)=u_v+\sum_X[e_{\{v\}}]F_X(U),\qquad
 q_v(X,Z;U)=[e_{\{v\}}]K_{X,Z}(U)
\end{equation}
has a finite elementary representation anchored at $o_v$, with
constant part $u_v$ and cost at most
\begin{equation}\label{eq:ope-single-price}
                         2a/w_v+2b/w_v.
\end{equation}
Indeed concatenate the finite lists just constructed. Adding lists
adds their costs, and multiplication by a scalar multiplies cost
by its modulus. The constant $u_v$ is recorded separately and costs
nothing in this representation; its later affordability is a
separate condition.

Each flagged elementary insertion is carried by its chosen index in
the countable activity dictionary. Adding the finitely many such
indexed functions gives a family $k_{v,j}(U)$ with a summable uniform
majorant. Summing it over all indices with flagged pair $(X,Z)$
gives exactly $q_v(X,Z;U)$. Thus a finite local representation also
provides the dominated indexed family required by the transfer
formula. The estimates concern the local representation and do not
require a canonically chosen index for each singleton.

\subsection{Exact separation of the higher terms}
Let $k_j^0(U)$ be the real background flagged activities and let
$E(U)$ be the ordinary background exponent. For a fixed configuration
$U$, denote its flag set by $\mathcal Z(U)$. For an array $L$ indexed
by flagged pairs set
\[
 \mathcal G_U(L)=
 \sum_{\substack{\mathcal P\text{ a set of flagged pairs}\\
           \text{supports pairwise disjoint},\;
           \bigcup_{(X,Z)\in\mathcal P}Z=\mathcal Z(U)}}
                 \prod_{(X,Z)\in\mathcal P}L_{X,Z}.
\]
The current density is
\begin{equation}\label{eq:ope-single-density}
 D(U)=e^{E(U)}\exp\left(\sum_vu_ve_{\{v\}}
                                     +\sum_XF_X(U)\right)
                         \mathcal G_U(K(U)).
\end{equation}
Assume the constant coefficient of $K_{X,Z}$ is the sum of
$k_j^0$ over dictionary indices with pair $(X,Z)$; assume also
the background activities are absolutely summable. These are the
vacuum consistency and validity conditions retained in the preceding
regrouping.

Let $F_X^{\ge2}$ and $K_{X,Z}^{\ge2}$ denote the arrays obtained by
discarding coefficients with fewer than two marks. Then, coefficient
by coefficient,
\begin{align*}
 \sum_vu_ve_{\{v\}}+\sum_XF_X
    &=\sum_vg_ve_{\{v\}}+\sum_XF_X^{\ge2},\\
 K_{X,Z}
    &=\sum_{j:\,(X_j,Z_j)=(X,Z)}k_j^0
       +\sum_v q_v(X,Z)e_{\{v\}}+K_{X,Z}^{\ge2}.
\end{align*}
Substitution in \eqref{eq:ope-single-density} proves the desired
local representation of the same density, with fields $g_v$,
flagged profiles $q_v$, and interactions of order at least two.
Discarding coefficients cannot increase a sum of nonnegative
coefficient majorants. Hence the two higher arrays retain hit bounds
$a$ and $b$. Their support, eligibility and invariance properties
are retained as well, since each remaining coefficient is unchanged.

For completeness, this representation has the expected linear
density. Put $r(U)=e^{E(U)}\mathcal G_U(k^0)$, the unmarked density.
Expanding a product to first order gives
\begin{equation}\label{eq:ope-single-coin}
 [e_{\{v\}}]D(U)=r(U)g_v(U)
 +e^{E(U)}\sum_{S\text{ admissible}}
          \sum_{j\in S}k_{v,j}(U)
                         \prod_{i\in S\setminus\{j\}}k_i^0(U).
\end{equation}
Here admissibility means that the indexed supports are disjoint
and their flags cover $\mathcal Z(U)$. The exponential contributes
its singleton coefficient $g_v$, while the flagged product contributes
exactly one differentiated factor. Replacing sums over flagged pairs
by sums over dictionary indices is legitimate: if
$B_0=\sum_j|k_j^0(U)|$ and $B_v=\sum_j|k_{v,j}(U)|$, the absolute
sum of all differentiated products, even without admissibility,
is bounded by $B_v\prod_j(1+|k_j^0(U)|)\le B_ve^{B_0}$.
The corresponding vacuum sum is bounded by $e^{B_0}$. Finite
expansions followed by absolute convergence establish the identity.

Equation~\eqref{eq:ope-single-coin} also explains why the ordinary
singleton is determined by its density and flagged profile whenever
$r(U)\ne0$: subtract the flagged term and divide by $r(U)$.
This is the uniqueness used when a transfer step first identifies
the singleton density rather than the ordinary field itself.

\subsection{Returning to the scale induction}
Apply this construction to the regrouped arrays of the preceding
section. If both hit bounds are at most $Z$, each grouped field has
cost at most $4Z/w_v$, and the higher interactions have hit bound
$Z$. The regrouping creates no separately recorded linear constant,
so here $u_v=0$. In particular, when $Z=1$, reducing the mark weight
from $w_v$ to $w'_v=w_v/8$ leaves
\[
             \text{singleton cost}\le\frac4{w_v}
                                  =\frac1{2w'_v}.
\]
The coefficient prices of the higher interactions only decrease
under this reduction. This proves the representation and the
singleton budget needed after a grouping operation, conditional on
the stated dictionary and background-validity properties.

Assume the input arrays and constants have the covariance inherited
from the preceding regrouping. Translations and coordinate symmetries
then preserve the ordinary
singletons by reindexing the finite support sum in
\eqref{eq:ope-single-pair}; conjugating symmetries also conjugate
its scalar coefficients. The flagged profiles transform in the
same way. No equivariance of the auxiliary choice of dictionary
indices is needed, because only their summed profiles occur in
the representation and in \eqref{eq:ope-single-coin}. This argument
applies to any finite number of slots in each copied cell.
The remaining analytic step is to replace these local fields by
degree-bounded coefficient fields while controlling the error
under transfer to the next scale. The singleton extraction alone
does not provide that projection estimate.
